\vspace{-5.0mm}
\section{Conclusion}
\vspace{-4mm}

In this work, we propose Equiformer, a graph neural network (GNN) combining the strengths of Transformers and equivariant features based on irreducible representations (irreps).
With irreps features, we build upon existing generic GNNs and Transformer networks by incorporating equivariant operations like tensor products.
We further propose equivariant graph attention, which incorporates multi-layer perceptron attention and non-linear messages.
Experiments on QM9, MD17 and OC20 demonstrate the effectiveness of Equiformer and ablation studies show the improvement of the proposed equivariant graph attention over typical attention in Transformers.

\begin{comment}
The broader impact lies in two aspects.
First, Equiformer demonstrates the possibility of adapting Transformers to domains such as physics and chemistry, where data can be represented as 3D atomistic graphs.
Second, Equiformer achieves more accurate approximations of quantum properties calculation.
%This has the potential to speed up molecular simulations and aid chemical applications such as discovering and optimizing catalysts~\cite{oc20}.
%However, these methods accelerate the ability to model quantum mechanical properties regardless of the intent of how the results will be used, and there is always a risk that there will be used for nefarious purposes.
We believe there is much more to be gained by harnessing these abilities for productive investigation of molecules and materials relevant to application such as energy, electronics, and pharmaceuticals~\cite{oc20}, than to be lost by applying these methods for adversarial purposes like creating hazardous chemicals.
Additionally, there are still substantial hurdles to go from the identification of a useful or harmful molecule to its large-scale deployment.

%However, at the same time, this could be of negative usage like creating hazardous chemicals, which is not the purpose of this work.
%Whether the proposed method has positive impacts or not would depend on datasets.
%The negative usage is not the purpose of this work.

\end{comment}


\section{Ethics Statement}
Equiformer achieves more accurate approximations of quantum properties calculation.
We believe there is much more to be gained by harnessing these abilities for productive investigation of molecules and materials relevant to application such as energy, electronics, and pharmaceuticals, than to be lost by applying these methods for adversarial purposes like creating hazardous chemicals.
Additionally, there are still substantial hurdles to go from the identification of a useful or harmful molecule to its large-scale deployment.

Moreover, we discuss several limitations of Equiformer and the proposed equivariant graph attention in Sec.~\ref{appendix:sec:limitations} in appendix.


\section{Reproducibility Statement }

We include details on architectures, hyper-parameters and training time in Sec.~\ref{appendix:subsec:qm9_training_details} (QM9), Sec.~\ref{appendix:subsec:md17_training_details} (MD17) and Sec.~\ref{appendix:subsec:oc20_training_details} (OC20).

%We submit our code reproducing the result of Equiformer on the task of $\alpha$ of QM9 dataset.
%Following the author guide, we will make a comment directed to the reviewers and area chairs and put a link to an anonymous repository.
The code for reproducing the results of Equiformer on QM9, MD17 and OC20 datasets is available at \texttt{\href{https://github.com/atomicarchitects/equiformer}{\color{crimson}{https://github.com/atomicarchitects/equiformer}}}.


%\vspace{-2mm}
\section*{Acknowledgement}
%\vspace{-2mm}
%We thank \texttt{e3nn}~\cite{e3nn} developers for the library and detailed guidelines and also thank Simon Batzner, Albert Musaelian, Mario Geiger, Johannes Brandstetter, and Rob Hesselink for their help in OC20 dataset.
We thank Simon Batzner, Albert Musaelian, Mario Geiger, Johannes Brandstetter, and Rob Hesselink for helpful discussions including help with the OC20 dataset.
We also thank the \texttt{e3nn}~\citep{e3nn} developers and community for the library and detailed documentation. 
We acknowledge the MIT SuperCloud and Lincoln Laboratory Supercomputing Center~\citep{supercloud} for providing high performance computing and consultation resources that have contributed to the research results reported within this paper.

Yi-Lun Liao and Tess Smidt were supported by DOE ICDI grant DE-SC0022215.